
\subsubsection{NLI Domain Mismatch}

The H-M2 result establishes that microsoft/deberta-large-mnli does not produce meaningful semantic aggregation on HaluEval-QA short factual QA responses under the lorenzkuhn/semantic_uncertainty configuration with N=5 stochastic samples. In 72.8\% of examples, all five stochastic responses are assigned to distinct semantic clusters, collapsing semantic entropy to a constant maximum-entropy value.

A plausible explanation is NLI domain mismatch: deberta-large-mnli was trained on MNLI, which consists of sentence pairs from news and fiction text, and was validated by Kuhn et al. (2023) on TriviaQA and Natural Questions — benchmarks whose responses are typically longer and more syntactically diverse than HaluEval-QA's 1–2 sentence factual answers. For short factual answers where the same semantic content is expressed in surface-variant forms (e.g., "The United States," "America," "the US"), the NLI model may classify surface-distinct responses as semantically non-equivalent, because its entailment threshold is calibrated for longer and more syntactically structured sentence pairs. This explanation is consistent with the data but has not been directly verified in the present work; ruling out alternative explanations (e.g., genuine semantic diversity in LLaMA-2-7B-chat responses at temperature = 1.0) would require testing with a QA-specific NLI model or examining individual response pairs.

The practical implication of the H-M2 finding is that semantic entropy's non-degeneracy cannot be assumed when applying the method to a new benchmark. The aggregation rate is a measurable precondition for non-degenerate semantic entropy, and reporting it alongside AUROC results would allow others to distinguish genuine discrimination from degenerate constant behavior.

\subsubsection{SelfCheckGPT Polarity Inversion Hypothesis}

The SelfCheckGPT-BERTScore result (AUROC = 0.3562, CI entirely below 0.5) indicates that BERTScore consistency is positively correlated with HaluEval-QA hallucination labels — the opposite of SelfCheckGPT's design assumption. One explanation, consistent with the benchmark's construction methodology, is a label polarity inversion: HaluEval-QA hallucinations are ChatGPT-generated confident-sounding incorrect answers. If LLaMA-2-7B-chat generates consistently wrong answers for questions with hallucinated labels — because the wrong answer is the model's preferred factual response — then BERTScore consistency would be higher for hallucinated examples than for factual examples, inverting SelfCheckGPT's expected correlation direction. Figure 5.4 shows that cluster count distributions are similar across hallucinated and factual labels, which is consistent with the model behaving similarly regardless of ground-truth label.

This hypothesis is not verified in the present work. Direct verification would require stratifying the per-example BERTScore consistency scores by hallucination label using the existing H-E1 outputs, which involves no new LLM inference. Alternative explanations — including BERTScore insensitivity to short factual QA responses and LLaMA-2-7B-chat-specific generation patterns — have not been ruled out.

\subsubsection{Token Entropy as a Null Baseline}

Token entropy (AUROC = 0.4829, CI [0.4585, 0.5090]) is consistent with zero discrimination on HaluEval-QA. This result indicates that LLaMA-2-7B-chat's per-token uncertainty does not systematically covary with HaluEval-QA hallucination labels under greedy inference. As with the other methods, this finding is specific to the LLaMA-2-7B-chat model and the HaluEval-QA benchmark; it should not be interpreted as a general failure of token entropy for hallucination detection.

\subsubsection{Limitations}

\textbf{Single model.} All executed experiments use LLaMA-2-7B-chat. The planned cross-model experiment (H-M3 with Mistral-7B-Instruct) was not executed. All findings are specific to LLaMA-2-7B-chat on HaluEval-QA; generalization to other model families has not been tested. The NLI aggregation failure is a property of deberta-large-mnli's behavior on HaluEval-QA response style; whether it would also occur with Mistral-7B-Instruct responses is unknown.

\textbf{ChatGPT-generated labels.} HaluEval-QA labels are generated by ChatGPT rather than human annotators. All three methods are subject to the same label characteristics, preserving between-method comparison validity. The near-random AUROC values for all methods are consistent with genuine method failure on this benchmark, or with label construction artifacts that suppress UQ signal correlation — distinguishing these would require applying the same pipeline to human-annotated benchmarks such as TriviaQA or Natural Questions with the same model.

\textbf{N=5 stochastic samples.} The NLI aggregation failure (72.8\% of examples at maximum cluster count, CI entirely below 0.30) is too large in magnitude to be plausibly explained by low N alone. However, whether increasing N (e.g., to N=20) would improve the aggregation rate has not been tested. Similarly, the SelfCheckGPT AUROC magnitude (AUROC = 0.3562, 0.144 below random) suggests a systematic effect unlikely to disappear with higher N, but this has not been verified.

\textbf{Unverified polarity inversion.} The proposed explanation for the SelfCheckGPT below-random result is a hypothesis, not an experimentally verified mechanism. Alternative explanations for AUROC = 0.3562 have not been systematically ruled out.

